% !TeX root = ../../main.tex

\section{Introducción}

Apache Hadoop es una plataforma creada para almacenar y procesar grandes
cantidades de datos mediante grupos de computadoras. Su funcionamiento
distribuido permite dividir el almacenamiento y el procesamiento entre varios
equipos, conocidos como nodos.

\section{¿Qué es Apache Hadoop?}

Apache Hadoop es un framework de código abierto utilizado para almacenar y
procesar grandes volúmenes de datos de manera distribuida. En lugar de depender
de una sola computadora, Hadoop utiliza un conjunto de máquinas conectadas que
forman un clúster.

Hadoop puede trabajar con datos estructurados, semiestructurados y no
estructurados. También permite ejecutar tareas en paralelo, lo que facilita
el procesamiento de conjuntos de datos de gran tamaño.

\section{Historia de Apache Hadoop}

El origen de Hadoop está relacionado con Nutch, un motor de búsqueda de código
abierto desarrollado por Doug Cutting y Mike Cafarella. Conforme aumentó la
cantidad de información de Internet, fue necesario encontrar una manera de
almacenar y procesar los datos entre varias computadoras.

Su desarrollo tomó como inspiración dos trabajos publicados por Google: Google
File System, en 2003, y MapReduce, en 2004. Estas ideas permitían distribuir
el almacenamiento y el procesamiento de datos.

En 2005 se implementaron MapReduce y un sistema de archivos distribuido en
Nutch. Posteriormente, Doug Cutting se incorporó a Yahoo!, donde continuó el
desarrollo del proyecto. Hadoop se convirtió en un proyecto de Apache y fue
adoptado por diferentes empresas. Hadoop 1.0 apareció en 2011, Hadoop 2 en
2013 y Hadoop 3.0 en 2017.

\section{Características principales}

Las características principales de Apache Hadoop son:

\begin{itemize}
    \item \textbf{Procesamiento distribuido:} divide una tarea entre varios
    nodos del clúster.

    \item \textbf{Escalabilidad:} permite agregar nuevos nodos cuando aumenta
    la cantidad de datos.

    \item \textbf{Tolerancia a fallos:} mantiene copias de los bloques de datos
    para evitar pérdidas cuando un nodo deja de funcionar.

    \item \textbf{Procesamiento paralelo:} varios nodos pueden trabajar al
    mismo tiempo sobre diferentes partes de un conjunto de datos.

    \item \textbf{Uso de hardware convencional:} puede utilizar computadoras
    de costo moderado en lugar de depender exclusivamente de servidores
    especializados.

    \item \textbf{Flexibilidad:} permite almacenar diferentes formatos de
    información.

    \item \textbf{Alta disponibilidad:} los datos replicados pueden recuperarse
    desde otros nodos.

    \item \textbf{Código abierto:} puede utilizarse, estudiarse y modificarse
    libremente.
\end{itemize}

\section{Módulos principales}

Apache Hadoop contiene cuatro módulos principales:

\begin{enumerate}
    \item \textbf{Hadoop Common:} contiene las bibliotecas y herramientas
    generales que necesitan los demás módulos.

    \item \textbf{HDFS:} es el sistema de archivos distribuido de Hadoop.
    Divide los archivos en bloques y los almacena en diferentes nodos.

    \item \textbf{YARN:} administra los recursos del clúster y programa la
    ejecución de las aplicaciones.

    \item \textbf{MapReduce:} es el modelo de procesamiento que divide los
    trabajos en tareas pequeñas y las ejecuta en paralelo.
\end{enumerate}

\section{Arquitectura de Apache Hadoop}

Hadoop utiliza una arquitectura distribuida de tipo maestro-trabajador. Sus
funciones principales se encuentran en las capas de almacenamiento y
procesamiento.

\subsection{Arquitectura de HDFS}

La arquitectura de HDFS está formada principalmente por:

\begin{itemize}
    \item \textbf{NameNode:} es el nodo principal. Conserva los metadatos del
    sistema, como los nombres de archivos, permisos y ubicación de los bloques.

    \item \textbf{DataNodes:} son los nodos trabajadores. Almacenan físicamente
    los bloques de datos y envían información sobre su estado al NameNode.

    \item \textbf{Cliente HDFS:} solicita operaciones para leer, escribir o
    eliminar archivos.
\end{itemize}

Los archivos se dividen en bloques que se distribuyen entre los DataNodes.
Cada bloque puede tener varias copias para mantener la información disponible
si ocurre una falla.

\subsection{Arquitectura de YARN}

YARN administra el procesamiento mediante los siguientes elementos:

\begin{itemize}
    \item \textbf{ResourceManager:} distribuye los recursos disponibles entre
    las aplicaciones.

    \item \textbf{NodeManager:} administra los recursos y tareas ejecutadas en
    cada nodo trabajador.

    \item \textbf{ApplicationMaster:} coordina la ejecución de una aplicación
    específica.

    \item \textbf{Contenedores:} proporcionan memoria y capacidad de
    procesamiento para ejecutar las tareas.
\end{itemize}

\section{¿Cómo trabaja Apache Hadoop?}

El funcionamiento básico de Hadoop puede resumirse en los siguientes pasos:

\begin{enumerate}
    \item Los datos se cargan en HDFS.
    \item HDFS divide los archivos en bloques.
    \item Los bloques se distribuyen y replican entre los DataNodes.
    \item El usuario envía un trabajo de procesamiento a YARN.
    \item YARN asigna los recursos y nodos necesarios para ejecutar el trabajo.
    \item MapReduce divide el trabajo en tareas Map.
    \item Las tareas Map procesan partes de los datos y generan resultados
    intermedios.
    \item Los resultados se ordenan y se envían a las tareas Reduce.
    \item Las tareas Reduce combinan los resultados.
    \item El resultado final se almacena nuevamente en HDFS.
\end{enumerate}

Hadoop procura ejecutar cada tarea cerca del nodo donde se encuentran los
datos. Esto disminuye la transferencia de información por la red.

\section{Ventajas}

Las principales ventajas de Apache Hadoop son:

\begin{itemize}
    \item Almacena y procesa grandes cantidades de datos.
    \item Puede aumentar su capacidad agregando más nodos.
    \item Distribuye las tareas para ejecutarlas en paralelo.
    \item Mantiene copias de los datos para enfrentar fallas.
    \item Puede utilizar hardware convencional.
    \item Trabaja con diferentes tipos y formatos de datos.
    \item Cuenta con una comunidad amplia y diversas herramientas relacionadas.
    \item Permite procesar los datos en el mismo nodo donde están almacenados.
\end{itemize}

\section{Desventajas}

Apache Hadoop también presenta algunas desventajas:

\begin{itemize}
    \item MapReduce se orienta principalmente al procesamiento por lotes, por
    lo que puede tener una latencia alta.

    \item No es la opción más adecuada para transacciones que requieren
    respuestas inmediatas.

    \item La instalación y administración de un clúster pueden ser complejas.

    \item Requiere conocimientos de sistemas distribuidos, redes y
    configuración de servidores.

    \item HDFS no trabaja de manera eficiente con una cantidad muy grande de
    archivos pequeños.

    \item La replicación aumenta el espacio de almacenamiento necesario.

    \item La seguridad y el monitoreo requieren configuración adicional.
\end{itemize}

\section{Tecnologías y productos relacionados}

El ecosistema de Hadoop contiene distintas herramientas que complementan sus
funciones:

\begin{itemize}
    \item \textbf{Apache Hive:} permite consultar datos mediante un lenguaje
    parecido a SQL.

    \item \textbf{Apache Pig:} utiliza el lenguaje Pig Latin para crear flujos
    de procesamiento de datos.

    \item \textbf{Apache HBase:} es una base de datos NoSQL que funciona sobre
    HDFS y permite acceder a registros específicos.

    \item \textbf{Apache Sqoop:} transfiere datos entre Hadoop y bases de datos
    relacionales.

    \item \textbf{Apache Flume:} recopila y transporta grandes cantidades de
    registros hacia Hadoop.

    \item \textbf{Apache Oozie:} programa y coordina flujos de trabajos.

    \item \textbf{Apache ZooKeeper:} ayuda a coordinar los servicios
    distribuidos.

    \item \textbf{Apache Spark:} proporciona procesamiento rápido en memoria,
    análisis, aprendizaje automático y procesamiento de flujos.

    \item \textbf{Apache Tez:} ofrece un motor alternativo para ejecutar
    procesos por lotes e interactivos.

    \item \textbf{Apache Ambari:} facilita la instalación, administración y
    supervisión de clústeres Hadoop.

    \item \textbf{Apache Avro:} proporciona un formato para serializar e
    intercambiar datos.
\end{itemize}

\section{Requerimientos de hardware y software}

Los requerimientos exactos dependen del volumen de información y del uso que
tendrá el clúster.

\subsection{Hardware}

Para una instalación de aprendizaje en una sola computadora se recomienda:

\begin{itemize}
    \item Procesador de 64 bits con dos o más núcleos.
    \item Al menos 4 GB de memoria RAM; 8 GB ofrecen una mejor experiencia.
    \item Entre 20 y 50 GB de almacenamiento disponible.
    \item Conexión de red cuando se utilicen varias máquinas.
\end{itemize}

En un ambiente de producción se necesitan varios servidores. Los DataNodes
deben contar con suficiente almacenamiento, memoria y capacidad de
procesamiento. También es importante utilizar una red rápida porque los nodos
intercambian información constantemente.

\subsection{Software}

Los requerimientos de software más comunes son:

\begin{itemize}
    \item Un sistema operativo Linux de 64 bits.
    \item Una versión del Java Development Kit compatible con la versión de
    Hadoop utilizada.
    \item Servicio SSH configurado para la comunicación entre los nodos.
    \item Una distribución de Apache Hadoop.
    \item Variables de entorno como \texttt{JAVA\_HOME} y
    \texttt{HADOOP\_HOME}.
    \item Nombres de host y direcciones de red configurados correctamente.
    \item Archivos de configuración para HDFS, YARN y MapReduce.
\end{itemize}

Para prácticas personales también puede ejecutarse Hadoop en una máquina
virtual o mediante contenedores.

\section{Implementación de Hadoop en la nube}

Hadoop puede implementarse en la nube utilizando máquinas virtuales o un
servicio administrado. Algunos ejemplos de servicios administrados son Amazon
EMR, Azure HDInsight y Google Cloud Dataproc.

Una implementación general puede realizarse con los siguientes pasos:

\begin{enumerate}
    \item Elegir el proveedor y la región donde se ejecutará el clúster.
    \item Seleccionar un servicio administrado o crear máquinas virtuales.
    \item Configurar la red, usuarios, permisos y reglas de seguridad.
    \item Seleccionar la cantidad y capacidad de los nodos.
    \item Crear el clúster e instalar los componentes necesarios.
    \item Configurar HDFS o conectarlo con el almacenamiento de objetos del
    proveedor.
    \item Transferir los datos mediante herramientas de ingesta.
    \item Ejecutar trabajos con MapReduce, Hive, Spark u otras herramientas.
    \item Supervisar el rendimiento, consumo y costo del clúster.
    \item Detener o reducir los recursos cuando ya no sean necesarios.
\end{enumerate}

Los servicios administrados simplifican la instalación, actualización y
escalamiento. Sin embargo, es necesario controlar los costos, los permisos de
acceso y la protección de los datos almacenados.

\section{Conclusión}

Apache Hadoop permite almacenar y procesar grandes cantidades de información
mediante un clúster de computadoras. HDFS se encarga del almacenamiento
distribuido, YARN administra los recursos y MapReduce ejecuta las tareas de
procesamiento. Aunque su administración puede ser compleja y no es la mejor
opción para todas las aplicaciones en tiempo real, continúa siendo una
tecnología importante para comprender el procesamiento distribuido y el
ecosistema de Big Data.